\documentclass[11pt,a4paper]{article}
\usepackage[T1]{fontenc}
\usepackage{lmodern,amsmath,amssymb,amsthm,mathtools,microtype}
\usepackage[margin=27mm]{geometry}
\usepackage[colorlinks=true,linkcolor=blue,citecolor=blue,urlcolor=blue]{hyperref}
\setlength{\emergencystretch}{2em}
\newtheorem{theorem}{Theorem}[section]
\newtheorem{proposition}[theorem]{Proposition}
\newtheorem{corollary}[theorem]{Corollary}
\newcommand{\Q}{\mathbb Q}\newcommand{\Z}{\mathbb Z}
\DeclareMathOperator{\Tr}{Tr}
\title{Finiteness from rational conditional detector products}
\author{}\date{17 November 2026}
\begin{document}\maketitle
\begin{abstract}
For the conditional elliptic detector, a rational field product
satisfying the pole implication has no nonexceptional odd contact.
An elementary valuation argument proves this assertion uniformly
in the auxiliary points. Faltings's theorem then restricts its rational
argument to a finite set. Requiring rational products in both coupled
blocks consequently permits only finitely many integer inputs.
The result leaves nonrational products and other detector designs open.
\end{abstract}

\section{The detector}

Let \(K=\Q(\sqrt2)\), \(D=11-6\sqrt2\), and
\[
 E:y^2=x^3-Dx,\qquad
 P=\left(\frac{57-28\sqrt2}{8},\frac{-660+417\sqrt2}{32}\right).
\]
Fix a finite set \(S\) of rational primes containing every prime
at most \(76176\). On the guarded domain of the conditional source
\cite{Conditional,Reweighted}, put
\[
 F(t)=\frac{t(t+1)(2t+1)(3t+1)(4t+1)(5t+2)}{120},\qquad
 w(t)=\frac{F(t)}{4(t-1)^8},
\]
\[
 W(t)=\prod_{i=1}^6(1+C_i^2/F(t)),\qquad
 V=\Tr(W^{-1}),\qquad b=4w/V^2.
\]
The mandatory tests include \(C_1=f(H)\) and
\(C_3=\overline{f(H-P)}\), where \(f=1/x\).
All original point, difference and reciprocal trace guards are retained.
We assume that the chosen pole predicate supplies
\begin{equation}\label{eq:pole}
 v_p(b)\text{ is not negative and odd for every }p\notin S.
\end{equation}
This implication is satisfied by each guarded native square test
in \cite{Native}; no converse is assumed here.

\section{Rational products and contacts}

\begin{proposition}\label{prop:local}
If an actual guarded source tuple satisfies \eqref{eq:pole} and
\(W\in\Q\), then \(v_p(F(t))\) is even for every \(p\notin S\).
\end{proposition}
\begin{proof}
If \(v_p(t)<0\), then \(v_p(F(t))=6v_p(t)\) is even.
If \(t\equiv1\pmod p\), then \(F(t)\equiv7\pmod p\) is a unit.
Every hypothetical odd valuation is therefore
\[
 \nu=v_p(F(t))=v_p(w)>0.
\]
Since \(W\) is rational, \(V=2/W\) and \(b=wW^2\).
Condition \eqref{eq:pole} implies
\[
 \nu+2v_p(W)\ge1,\qquad v_p(W)>-\nu/2.
\]
At either place above \(p\), every factor has valuation
\[
 \min(0,2v(C_i)-\nu)\le0.
\]
There is no cancellation of the two terms, because their valuations
have different parity. The sum of these valuations is the same
\(v_p(W)\) at both places. Each factor valuation is at least that
sum, so \(v(C_i)>\nu/4>0\) at both places for every \(i\).

At one fixed place, positivity for \(C_1\), together with the
conjugate-place assertion for \(C_3\), puts both \(H\) and
\(H-P\) in the formal group. Their difference \(P\) would also
reduce to the identity. This is impossible: outside \(S\) the
curve has good reduction and \(P\) is an affine integral point.
The argument applies to split and inert primes alike.
\end{proof}

\begin{corollary}\label{cor:finite}
The possible rational arguments of tuples in Proposition~\ref{prop:local}
belong to a finite set depending only on \(S\).
\end{corollary}
\begin{proof}
Their nonzero \(F(t)\)'s belong to the finite group of rational
square classes generated by \(-1\) and the primes of \(S\).
For each representative \(d\), the equation \(F(t)=d y^2\)
has a smooth projective model of genus two: the degree-six polynomial
has six distinct roots. Faltings's theorem \cite{Faltings} gives
finitely many rational points on each model. Their finite union
gives the asserted finite set of rational arguments.
\end{proof}

\begin{theorem}\label{thm:two}
Fix an odd \(B\ge3\) and a positive even integer \(e\).
In the coupled conditional source
\[
 t_0=\lambda,\qquad t_1=a_1\lambda,\qquad
 a_1=(X^2+B)^e+\Tr(x(A))^e,
\]
only finitely many integer inputs \(X\) admit guarded tuples with
both products \(W(t_0),W(t_1)\in\Q\) satisfying \eqref{eq:pole}.
\end{theorem}
\begin{proof}
Let \(\mathcal T\) be the finite set in Corollary~\ref{cor:finite}.
The slope guard gives \(t_0\ne0\). Hence
\[
 a_1\in\{t'/t:t,t'\in\mathcal T,\ t\ne0\},
\]
a fixed finite set. But \(a_1\ge(X^2+B)^e\), which is unbounded
on integer inputs. Thus \(|X|\) is bounded.
\end{proof}

The local proposition needs neither fixed auxiliary points nor any
prescribed multiplication relation between them. Requiring a rational
product in just one block restricts that block's argument to a finite
set; the theorem does not exclude every such mixed construction.
The proof concerns the displayed conditional weights. The additional
normalizers and scalar factors of the general source require a
separate argument. No conclusion about all rational witnesses on
the native square covers follows from this theorem.

\begin{thebibliography}{9}
\bibitem{Faltings} G. Faltings, \emph{Endlichkeitss\"atze f\"ur abelsche
Variet\"aten \"uber Zahlk\"orpern}, Invent. Math. \textbf{73} (1983),
349--366. \href{https://doi.org/10.1007/BF01388432}{doi:10.1007/BF01388432}.
\bibitem{Conditional} \emph{Thirteen-variable definitions from direct
input-height amplification}, repository manuscript, 10 November 2026.
\bibitem{Reweighted} \emph{Reweighted elliptic detectors}, repository
manuscript, 16 November 2026.
\bibitem{Native} \emph{A guarded rational-square criterion for normalized
elliptic detectors}, repository manuscript, 16 November 2026.
\end{thebibliography}
\end{document}
